%% defense_questions.tex
%% Thesis Defense Preparation: Potential Questions & Answer Strategies
%% LBEADS-NET -- Saad Zubairi, NYU Tandon, Advisor: Prof. Ivan Selesnick
%% Generated 2026-04-23
%% ============================================================================

\documentclass[11pt, letterpaper]{article}

% ── Geometry ─────────────────────────────────────────────────────────────────
\usepackage[margin=1in]{geometry}

% ── Fonts & encoding ────────────────────────────────────────────────────────
\usepackage[T1]{fontenc}
\usepackage{lmodern}
\usepackage{microtype}

% ── Colours & boxes ─────────────────────────────────────────────────────────
\usepackage[dvipsnames,svgnames,x11names]{xcolor}
\usepackage[most]{tcolorbox}

% ── Lists ───────────────────────────────────────────────────────────────────
\usepackage{enumitem}

% ── Hyperlinks & ToC ───────────────────────────────────────────────────────
\usepackage[bookmarks,hidelinks,
            colorlinks=true,
            linkcolor=NavyBlue,
            urlcolor=RoyalBlue]{hyperref}

% ── Math ────────────────────────────────────────────────────────────────────
\usepackage{amsmath,amssymb}

% ── Misc ────────────────────────────────────────────────────────────────────
\usepackage{booktabs}
\usepackage{titlesec}

% ============================================================================
% Colour palette
% ============================================================================
\definecolor{QBoxFrame}{HTML}{2B6CB0}   % steel-blue frame for questions
\definecolor{QBoxBG}{HTML}{EBF4FA}      % very light blue fill
\definecolor{ABoxFrame}{HTML}{276749}   % forest-green frame for answers
\definecolor{ABoxBG}{HTML}{F0FFF4}      % very light green fill
\definecolor{WarnFrame}{HTML}{C53030}   % red frame for warnings
\definecolor{WarnBG}{HTML}{FFF5F5}      % very light red fill
\definecolor{SectionBlue}{HTML}{1A365D} % dark navy for section headings

% ============================================================================
% tcolorbox styles
% ============================================================================
\newtcolorbox{questionbox}[1][]{%
  enhanced,
  colback=QBoxBG,
  colframe=QBoxFrame,
  coltitle=white,
  fonttitle=\bfseries\sffamily,
  left=4pt, right=4pt, top=4pt, bottom=4pt,
  boxrule=1pt,
  attach boxed title to top left={yshift=-2mm, xshift=4mm},
  boxed title style={colback=QBoxFrame, colframe=QBoxFrame,
                     sharp corners, size=small},
  sharp corners=south,
  title={#1}
}

\newtcolorbox{answerbox}{%
  enhanced,
  colback=ABoxBG,
  colframe=ABoxFrame,
  left=4pt, right=4pt, top=4pt, bottom=4pt,
  boxrule=0.8pt,
  fontupper=\small,
  sharp corners=south,
  title={\sffamily\bfseries Answer Strategy},
  coltitle=white,
  attach boxed title to top left={yshift=-2mm, xshift=4mm},
  boxed title style={colback=ABoxFrame, colframe=ABoxFrame,
                     sharp corners, size=small}
}

\newtcolorbox{warnbox}[1][]{%
  enhanced,
  colback=WarnBG,
  colframe=WarnFrame,
  coltitle=white,
  fonttitle=\bfseries\sffamily,
  left=4pt, right=4pt, top=4pt, bottom=4pt,
  boxrule=1pt,
  attach boxed title to top left={yshift=-2mm, xshift=4mm},
  boxed title style={colback=WarnFrame, colframe=WarnFrame,
                     sharp corners, size=small},
  sharp corners=south,
  title={#1}
}

% ============================================================================
% Section formatting
% ============================================================================
\titleformat{\section}
  {\Large\bfseries\sffamily\color{SectionBlue}}
  {\thesection}{1em}{}[\vspace{2pt}\titlerule]
\titleformat{\subsection}
  {\large\bfseries\sffamily\color{SectionBlue!80!black}}
  {\thesubsection}{1em}{}

% ============================================================================
% Question counter
% ============================================================================
\newcounter{qcounter}
\newcommand{\question}[1]{%
  \stepcounter{qcounter}%
  \begin{questionbox}[Q\theqcounter]%
    \itshape #1%
  \end{questionbox}%
}

% ============================================================================
% Title
% ============================================================================
\title{%
  \vspace{-1cm}
  {\sffamily\LARGE Thesis Defense Preparation}\\[6pt]
  {\sffamily\Large Potential Questions \& Answer Strategies}\\[12pt]
  {\normalfont\large
   \textbf{LBEADS-NET}: Algorithm Unrolling for Learned Baseline\\
   Estimation and Denoising in Chromatographic Signals}
}
\author{%
  Saad Zubairi\\
  \small MS Electrical Engineering, NYU Tandon School of Engineering\\
  \small Advisor: Professor Ivan Selesnick
}
\date{April 2026}

% ============================================================================
\begin{document}
\maketitle
\thispagestyle{empty}

\vspace{1em}
\begin{center}
\fbox{\parbox{0.85\textwidth}{\centering\small\sffamily
This document contains \textbf{65 anticipated defense questions} organized into
eleven categories, each with a suggested answer strategy and references to
specific thesis sections, tables, and equations.  Section~11 contains the
\textbf{killer questions}---the hardest, most uncomfortable challenges.
A final section catalogues known inconsistencies to fix.}}
\end{center}

\newpage
\tableofcontents
\newpage

% ============================================================================
%  1. FOUNDATIONAL / "WHY THIS PROBLEM"
% ============================================================================
\section{Foundational / ``Why This Problem'' Questions}

\setcounter{qcounter}{0}

% ── Q1 ──────────────────────────────────────────────────────────────────────
\question{Why chromatographic baseline correction specifically?  Why not
apply algorithm unrolling to general signal denoising, where the
impact would be broader?}

\begin{answerbox}
\textbf{Key points:}
\begin{itemize}[nosep]
  \item Chromatography has a \emph{domain-specific structure} that generic
        denoisers ignore: the signal decomposes into exactly three additive
        components---baseline ($\mathbf{b}$), peaks ($\mathbf{p}$), and
        noise ($\boldsymbol{\varepsilon}$)---with distinct spectral and
        sparsity priors (Sec.~2.1).
  \item BEADS already encodes those priors via asymmetric penalty and
        highpass filtering.  Unrolling preserves the interpretable
        structure while letting each layer adapt.
  \item Regulatory context matters: pharmaceutical QC (USP~$\langle 621\rangle$)
        and metabolomics both require \emph{auditable} preprocessing, which
        favours interpretable, model-based architectures over black-box nets.
  \item General denoising is well-served by DnCNN / U-Net variants;
        baseline estimation is an under-explored niche where
        model-based deep learning is a natural fit.
\end{itemize}
\textbf{Thesis references:} Sec.~1.1 (motivation), Sec.~2.1 (chromatographic model), Sec.~2.3.3 (BEADS formulation).
\end{answerbox}

% ── Q2 ──────────────────────────────────────────────────────────────────────
\question{What is wrong with just using arPLS or airPLS?  They are simpler,
they do not need training data, and they are already widely adopted
in chemometrics.  Convince me that the added complexity of unrolling
is worth it.}

\begin{answerbox}
\textbf{Key points:}
\begin{itemize}[nosep]
  \item arPLS/airPLS are iteratively-reweighted penalized least-squares
        methods.  They require manual selection of the smoothness
        parameter $\lambda$ and the asymmetry weight $p$, which are
        signal-dependent (Sec.~2.3.1--2.3.2).
  \item They perform \emph{only} baseline estimation---they do not
        simultaneously denoise.  BEADS and LBEADS-NET jointly estimate
        baseline and denoise.
  \item The core contribution is not ``better numbers than arPLS'' but
        \emph{demonstrating the unrolling paradigm} in this domain and
        exposing the fundamental $\ell_1$ leakage limitation.
  \item Concede honestly: on well-tuned, non-dense-peak data, arPLS
        will often suffice.  LBEADS-NET is interesting as a research
        vehicle for understanding leakage.
\end{itemize}
\textbf{Thesis references:} Table~2.1 (baseline methods comparison), Sec.~2.3.1--2.3.2 (arPLS, airPLS).
\end{answerbox}

% ── Q3 ──────────────────────────────────────────────────────────────────────
\question{Why does baseline drift matter quantitatively?  What error
magnitudes are we talking about in real instruments, and how does
that translate to, say, drug concentration accuracy?}

\begin{answerbox}
\textbf{Key points:}
\begin{itemize}[nosep]
  \item Column bleed, solvent gradients, and detector drift produce
        baseline offsets of 1--10\% of the tallest peak height in HPLC
        (Sec.~2.1.2).
  \item In pharmaceutical QC, area-percent quantitation is the norm.
        A 2\% baseline offset can shift a peak area by 5--15\%, which
        exceeds the typical $\pm$2\% USP acceptance criterion.
  \item Metabolomics (hundreds of analytes per run) amplifies the
        problem: even small per-peak errors compound across the profile.
  \item Acknowledge: we do not empirically quantify this in the thesis
        because we work with synthetic data.  Cite the chromatography
        literature for representative magnitudes.
\end{itemize}
\textbf{Thesis references:} Sec.~2.1.2 (baseline sources), Sec.~1.1 (quantitative motivation).
\end{answerbox}

% ── Q4 ──────────────────────────────────────────────────────────────────────
\question{Who actually uses this in practice---pharmaceutical QC,
metabolomics, environmental labs?  Have you talked to anyone who
processes real chromatograms and asked whether they need this?}

\begin{answerbox}
\textbf{Key points:}
\begin{itemize}[nosep]
  \item Be honest: this is a methods thesis, not a deployment study.
        No practitioner interviews were conducted.
  \item Frame the contribution as \emph{foundational analysis}: we
        unrolled BEADS, exposed the $\ell_1$ leakage limit, and
        documented the failure modes.  That knowledge benefits future
        applied work.
  \item Note that Professor Selesnick co-authored the original BEADS
        algorithm, so this work sits in a signal-processing research
        lineage rather than a product-development one.
  \item Mention potential users: Agilent OpenLab, Waters Empower,
        Shimadzu LabSolutions all rely on classical baseline algorithms
        that could benefit from learned parameter tuning.
\end{itemize}
\textbf{Thesis references:} Sec.~1.1, Sec.~6 (future work).
\end{answerbox}

% ── Q5 ──────────────────────────────────────────────────────────────────────
\question{Why not just use a black-box neural network---a U-Net or a
1-D ResNet?  They clearly generalise better and do not suffer from
$\ell_1$ leakage.  What does unrolling buy you over a purely
learned architecture?}

\begin{answerbox}
\textbf{Key points:}
\begin{itemize}[nosep]
  \item \textbf{Parameter efficiency:} LBEADS-NET has \textbf{100 learnable
        scalars} ($5 \times 20$ layers).  A comparable 1-D U-Net would
        have $10^5$--$10^6$ parameters, requiring far more training data
        than our 200 synthetic signals (Sec.~3.2, Sec.~4.1).
  \item \textbf{Interpretability:} every learned parameter has a physical
        meaning---$\lambda_0$ controls fidelity, $\lambda_1$/$\lambda_2$
        control sparsity, $r$ controls penalty shape, $\eta$ is the
        step size.  This matters for regulatory contexts.
  \item \textbf{Inductive bias:} the unrolled architecture \emph{guarantees}
        the three-component decomposition structure.  A black-box net
        must learn this from data alone.
  \item \textbf{Honest concession:} black-box nets would likely achieve
        lower reconstruction error given enough data.  The trade-off
        is interpretability and data efficiency versus raw performance.
\end{itemize}
\textbf{Thesis references:} Sec.~2.4 (algorithm unrolling theory), Sec.~3.2 (architecture), Table~3.2 (parameter count).
\end{answerbox}


% ============================================================================
%  2. ALGORITHM UNROLLING THEORY
% ============================================================================
\section{Algorithm Unrolling Theory}

% ── Q6 ──────────────────────────────────────────────────────────────────────
\question{What convergence guarantees does LBEADS-NET inherit from the
classical BEADS algorithm?  Can you prove that the learned iterates
converge to anything meaningful?}

\begin{answerbox}
\textbf{Key points (this is a trap question---acknowledge it):}
\begin{itemize}[nosep]
  \item Classical BEADS is derived via majorisation-minimisation (MM)
        and converges to the global optimum of its convex cost function
        (Sec.~2.3.3, Eq.~2.24).
  \item \textbf{Untied parameters break MM convergence.}  Each layer
        minimises a \emph{different} surrogate, so the sequence of
        iterates does not descend a single objective.
  \item This is a known trade-off in all algorithm unrolling:
        we sacrifice convergence guarantees for task-specific performance
        via end-to-end training (Sec.~2.4.3).
  \item Cite Chen \& Pock (2017) and Monga et al.\ (2021): unrolled
        networks should be understood as \emph{feed-forward architectures
        with structured layers}, not as truncated optimisers.
  \item The network is still well-posed in that training converges
        and test loss is stable---but ``convergence of the iterates''
        in the MM sense does not apply.
\end{itemize}
\textbf{Thesis references:} Sec.~2.4.3 (unrolling trade-offs), Sec.~3.3 (ISTA reformulation).
\end{answerbox}

% ── Q7 ──────────────────────────────────────────────────────────────────────
\question{Why ISTA over FISTA?  You lose the $O(1/k^2)$ momentum
acceleration for free.  Was this a deliberate choice or just simpler
to implement?}

\begin{answerbox}
\textbf{Key points:}
\begin{itemize}[nosep]
  \item The primary reason was \textbf{avoiding gradient vanishing through
        $K \times K_{\text{inner}}$ depth}.  Classical BEADS uses a
        conjugate-gradient inner loop; unrolling both outer MM and inner
        CG iterations creates excessive depth (Sec.~3.3).
  \item ISTA replaces the inner CG loop entirely: each unrolled layer
        is a single proximal-gradient step, giving depth $K$ instead
        of $K \times K_{\text{inner}}$.
  \item FISTA's momentum term introduces a non-monotone trajectory
        that complicates backpropagation through the unrolled graph
        (the Nesterov correction depends on the \emph{previous two}
        iterates, increasing memory footprint).
  \item In the unrolled setting with untied parameters, the step-size
        $\eta^{(k)}$ is learned per layer and effectively subsumes much
        of the acceleration benefit.
  \item Concede: FISTA layers are a natural extension for future work.
\end{itemize}
\textbf{Thesis references:} Sec.~3.3 (ISTA reformulation), Sec.~2.4.2 (LISTA and ISTA-based unrolling).
\end{answerbox}

% ── Q8 ──────────────────────────────────────────────────────────────────────
\question{How does your approach differ from LISTA fundamentally?  LISTA
learns entire weight matrices, whereas you learn only scalars.  Is
LBEADS-NET actually ``unrolling'' or just ``parameterised iteration''?}

\begin{answerbox}
\textbf{Key points:}
\begin{itemize}[nosep]
  \item LISTA (Gregor \& LeCun, 2010) unrolls ISTA for sparse coding:
        it learns the dictionary $\mathbf{W}$ and the Gram matrix
        $\mathbf{S}$ as free matrices, so each layer is a full affine
        transform (Sec.~2.4.2).
  \item LBEADS-NET retains the \emph{fixed structural operators}
        ($\mathbf{H}$ highpass filter, $\mathbf{D}$ difference matrices)
        and only untie the \emph{scalar regularisation weights} across
        layers---5 scalars per layer.
  \item This is a design choice for interpretability and parameter
        efficiency, not a limitation.  The structural operators encode
        domain knowledge (filter design, finite differences) that should
        \emph{not} be made freely learnable with only 200 training signals.
  \item It is still correctly termed ``unrolling'': each layer corresponds
        to one iteration of a known algorithm with layer-specific parameters.
        The Monga et al.\ (2021) taxonomy calls this
        ``algorithm unrolling with untied hyperparameters.''
\end{itemize}
\textbf{Thesis references:} Sec.~2.4.2 (LISTA), Sec.~3.2 (LBEADS-NET architecture), Table~3.2.
\end{answerbox}

% ── Q9 ──────────────────────────────────────────────────────────────────────
\question{What happens if you tie all parameters across layers---that is,
force $\lambda_0^{(k)} = \lambda_0$ for all $k$?  Does it reduce to
classical BEADS?  Have you tried it?}

\begin{answerbox}
\textbf{Key points:}
\begin{itemize}[nosep]
  \item With tied parameters, each layer applies the \emph{same}
        proximal-gradient step, so the unrolled network becomes
        $K$ iterations of ISTA applied to the BEADS cost function
        with \emph{learned} but \emph{shared} regularisation weights.
  \item It does \textbf{not} exactly reduce to classical BEADS because
        (a) BEADS uses MM with CG inner solve, not ISTA, and (b) the
        $K=20$ iteration budget may not suffice for convergence of the
        tied-parameter ISTA.
  \item The tied-parameter variant is an important ablation that
        isolates the benefit of \emph{per-layer specialisation}.
  \item Concede: this ablation is listed in Table~5.3 but has not
        been run yet.  It is planned before the final submission.
  \item The emergent parameter specialisation (Sec.~5.2.3) suggests
        that untying is beneficial: early layers learn coarse
        separation, later layers learn aggressive refinement.
\end{itemize}
\textbf{Thesis references:} Sec.~3.2 (architecture), Sec.~5.2.3 (parameter analysis), Table~5.3 (ablation, currently empty).
\end{answerbox}

% ── Q10 ─────────────────────────────────────────────────────────────────────
\question{The BEADS cost function is convex.  Does unrolling preserve
convexity of the end-to-end mapping?  If not, what are the
optimisation implications for training?}

\begin{answerbox}
\textbf{Key points:}
\begin{itemize}[nosep]
  \item \textbf{No.}  The end-to-end mapping
        $\mathbf{y}\mapsto(\hat{\mathbf{b}},\hat{\mathbf{p}})$ is a
        composition of 20 layers, each involving soft-thresholding
        (non-linear) and matrix multiplications parameterised by
        learned scalars.  The composition is \emph{not} convex in the
        parameters $\{\lambda_0^{(k)},\ldots,\eta^{(k)}\}$.
  \item The \emph{training loss landscape} is therefore non-convex,
        which is standard for neural networks.
  \item However, each \emph{individual layer} is a proximal operator
        applied to a convex penalty, which provides local regularity
        (Lipschitz continuity of the proximal map).
  \item Practically, the 3-stage curriculum (Sec.~3.8) and progressive
        loss introduction help navigate the non-convex landscape by
        providing a good initialisation path.
\end{itemize}
\textbf{Thesis references:} Sec.~2.3.3 (BEADS convexity), Sec.~3.3 (ISTA layers), Sec.~3.8 (curriculum training).
\end{answerbox}


% ============================================================================
%  3. ARCHITECTURE & DESIGN CHOICES
% ============================================================================
\section{Architecture \& Design Choices}

% ── Q11 ─────────────────────────────────────────────────────────────────────
\question{Why $K{=}20$ layers?  Your ablation table (Table~5.3) is
completely empty.  How do you justify this choice without any
empirical evidence?}

\begin{answerbox}
\textbf{Key points:}
\begin{itemize}[nosep]
  \item Honest answer: $K{=}20$ was chosen heuristically based on
        (a)~the classical BEADS default of 20--30 outer iterations
        and (b)~GPU memory constraints with $N{=}4096$ and dense matrices.
  \item The parameter analysis in Sec.~5.2.3 was conducted with $K{=}6$
        and $N{=}1024$, which is not directly comparable.
  \item Acknowledge that the empty ablation table is a gap.
        State that $K \in \{6, 10, 15, 20, 30\}$ should be swept.
  \item Qualitative argument: BEADS typically converges in
        $\sim$15--20 iterations, so the unrolled variant should need
        \emph{at most} that many layers (and likely fewer, since
        parameters are untied).
  \item Cite the LISTA literature: Borgerding et al.\ (2017) show
        that unrolled ISTA with 15--20 layers matches 100+ iterations
        of classical ISTA.
\end{itemize}
\textbf{Thesis references:} Table~5.3 (empty ablation), Sec.~5.2.3 (6-layer analysis), Sec.~3.2.
\end{answerbox}

% ── Q12 ─────────────────────────────────────────────────────────────────────
\question{Your implementation uses dense $N \times N$ matrices with
$O(N^2)$ memory.  For $N{=}4096$ in float64, each matrix is 128\,MB.
With 20 layers, you are storing gigabytes of intermediate buffers.
Why not exploit the banded structure of $\mathbf{H}$ and
$\mathbf{D}$?}

\begin{answerbox}
\textbf{Key points:}
\begin{itemize}[nosep]
  \item You are correct---this is the single largest engineering
        limitation of the current implementation.
  \item The highpass filter $\mathbf{H}$ is a banded Toeplitz matrix
        and $\mathbf{D}$ (first/second difference) is banded with
        bandwidth 1--2.  Products $\mathbf{H}^T\mathbf{H}$ and
        $\mathbf{D}^T\mathbf{D}$ are also banded.
  \item The reason for dense matrices: the filter $\mathbf{H}$ is
        constructed via \texttt{scipy.signal} and converted to a dense
        NumPy array before wrapping in a PyTorch tensor.  Implementing
        banded-matrix operations natively in PyTorch would require
        custom CUDA kernels or use of \texttt{torch.sparse}.
  \item This is an engineering priority for future work.  Banded storage
        would reduce memory from $O(N^2)$ to $O(N \cdot b)$ where $b$
        is the bandwidth ($b \ll N$), enabling $N{=}65536$+ signals.
  \item Practically, the current implementation fits on a single GPU
        for $N{=}4096$ only because of small batch sizes.
\end{itemize}
\textbf{Thesis references:} Sec.~3.4 (filter construction), Sec.~6 (future work).
\end{answerbox}

% ── Q13 ─────────────────────────────────────────────────────────────────────
\question{You say $f_c$ (the highpass cutoff frequency) is not learnable
because of the ``SciPy--PyTorch autograd boundary.''  But why not
implement the entire filter construction in PyTorch?  The
Butterworth poles-and-zeros are closed-form.}

\begin{answerbox}
\textbf{Key points:}
\begin{itemize}[nosep]
  \item The current implementation calls
        \texttt{scipy.signal.butter()} and
        \texttt{scipy.signal.tf2zpk()}, which are not differentiable
        through PyTorch's autograd (Sec.~3.4).
  \item In principle, the bilinear-transform design of a Butterworth
        filter \emph{can} be written as closed-form operations on
        $f_c$ (pole placement on the unit circle), and these operations
        are differentiable.
  \item The difficulty is constructing the \emph{banded matrix}
        $\mathbf{H}(f_c)$ from those poles in a way that is compatible
        with efficient batched backpropagation.
  \item A more practical route: parameterise the filter as a
        learnable FIR filter of fixed order and initialise it from the
        Butterworth design.  This avoids the IIR stability issues.
  \item This is listed as future work (Sec.~6) and would add one more
        learnable parameter per layer (101 total scalars).
\end{itemize}
\textbf{Thesis references:} Sec.~3.4 (filter construction), Eq.~3.5--3.6, Sec.~6 (future work).
\end{answerbox}

% ── Q14 ─────────────────────────────────────────────────────────────────────
\question{What is the effective receptive field of LBEADS-NET?  Each layer
operates on the full signal via dense matrix multiplication, so is
the receptive field trivially the entire signal?}

\begin{answerbox}
\textbf{Key points:}
\begin{itemize}[nosep]
  \item Yes---formally, every output sample depends on every input
        sample after a single layer, because $\mathbf{H}$ is an IIR
        filter realised as a dense matrix (Sec.~3.4).
  \item However, the \emph{effective} receptive field (in the
        Luo et al.\ 2016 sense) is narrower because the energy of
        $\mathbf{H}$ decays exponentially away from the diagonal.
  \item The practical receptive field is governed by the filter order
        and $f_c{=}0.006$: the $-3$\,dB point at $0.6\%$ of Nyquist
        corresponds to a spatial scale of $\sim$167 samples.
  \item For baseline estimation, a large receptive field is
        \emph{desirable}: the baseline is a slowly-varying component
        that depends on global signal context.
  \item For peak estimation, local structure matters more, which is
        why the sparsity penalties ($\ell_1$, TV) act element-wise
        or on adjacent differences.
\end{itemize}
\textbf{Thesis references:} Sec.~3.4 (filter design), Sec.~2.3.3 (BEADS formulation).
\end{answerbox}

% ── Q15 ─────────────────────────────────────────────────────────────────────
\question{Why float64 instead of float32?  Is this a fundamental numerical
requirement of the algorithm, or is it just an engineering choice
that doubles your memory usage?}

\begin{answerbox}
\textbf{Key points:}
\begin{itemize}[nosep]
  \item It is partially fundamental: the highpass filter $\mathbf{H}$
        at $f_c{=}0.006$ has eigenvalues spanning several orders of
        magnitude.  In float32, the matrix
        $\mathbf{H}^T\mathbf{H} + \lambda_0 \mathbf{I}$ can become
        poorly conditioned and produce NaN gradients during
        backpropagation.
  \item Empirically, early experiments in float32 showed training
        instability (gradient overflow in layers 15--20).
  \item However, with banded solvers or mixed-precision strategies
        (float64 for the filter, float32 for the rest), it should be
        possible to reduce memory.  This was not explored.
  \item The ISTA reformulation partially mitigates the conditioning
        issue (no matrix inverse required), so float32 may be viable
        with careful step-size initialisation.
\end{itemize}
\textbf{Thesis references:} Sec.~3.4 (numerical considerations), Sec.~4.2 (training configuration).
\end{answerbox}

% ── Q16 ─────────────────────────────────────────────────────────────────────
\question{You mention using softplus for non-negativity constraints in v7
but not in v5.  Softplus versus ReLU---what is the difference in
this context, and why does it matter?}

\begin{answerbox}
\textbf{Key points:}
\begin{itemize}[nosep]
  \item The regularisation parameters $\lambda_0, \lambda_1, \lambda_2,
        r, \eta$ must be strictly positive.  A reparameterisation is
        needed: the network stores unconstrained $\theta^{(k)}$ and
        maps to $\lambda^{(k)} = g(\theta^{(k)})$.
  \item \textbf{ReLU:} $g(\theta) = \max(0, \theta)$.  Has zero
        gradient for $\theta < 0$ (``dead parameter'' problem) and
        is non-differentiable at 0.
  \item \textbf{Softplus:} $g(\theta) = \log(1 + e^\theta)$.
        Smooth, strictly positive, non-zero gradient everywhere.
  \item In v7, softplus was adopted after observing that some
        parameters collapsed to zero under ReLU and could not recover
        (gradient was exactly zero).
  \item Softplus provides a mild bias toward small positive values
        rather than exact zeros, which is physically appropriate---the
        regularisation weights should not be zero.
\end{itemize}
\textbf{Thesis references:} Sec.~3.6 (parameter constraints), Sec.~3.7.1 (v7 modifications).
\end{answerbox}


% ============================================================================
%  4. LOSS FUNCTION DEEP DIVES
% ============================================================================
\section{Loss Function Deep Dives}

% ── Q17 ─────────────────────────────────────────────────────────────────────
\question{You have 12 loss terms with 12 hyperparameters (the weights).
Classical BEADS has 6 parameters.  Haven't you just replaced one
tuning problem with a larger one?  How is this progress?}

\begin{answerbox}
\textbf{Key points:}
\begin{itemize}[nosep]
  \item Fair criticism---acknowledge it directly.
  \item The critical difference: BEADS's 6 parameters must be tuned
        \emph{per signal} (or per instrument/method), whereas the
        12 loss weights are set \emph{once during training} and the
        network then adapts its 100 per-layer parameters to each new
        signal at inference time.
  \item The loss weights are meta-parameters that define the
        \emph{training objective}, not the inference-time model.
        At test time, LBEADS-NET has zero user-tunable parameters.
  \item Concede: the proliferation of loss terms (especially the
        12-term v7 variant) is itself a finding---Experiment 1.5 shows
        that the full 12-term loss \emph{hurts} performance due to
        competing gradients (Sec.~5.2).
  \item The thesis's \emph{negative result}---that more loss terms do
        not help---is itself a contribution.
\end{itemize}
\textbf{Thesis references:} Table~3.3 (loss terms), Sec.~5.2.1 (Exp.~1.5 results), Sec.~3.7.
\end{answerbox}

% ── Q18 ─────────────────────────────────────────────────────────────────────
\question{The orthogonality loss $\mathcal{L}_{\text{orth}}$ uses
\texttt{.detach()}'d ground-truth-derived weights.  Isn't this a
form of label leakage during training?  You are giving the network
access to the answer.}

\begin{answerbox}
\textbf{Key points:}
\begin{itemize}[nosep]
  \item Clarify what happens: the ground-truth peak signal is used to
        construct a \emph{weighting mask} that emphasises peak regions,
        and this mask is detached (no gradient flows through it).
        The loss penalises $\langle \hat{\mathbf{b}}, \mathbf{w} \rangle$
        where $\mathbf{w}$ is the detached mask.
  \item This is \textbf{not} label leakage in the test-time sense:
        the mask is used only during training (supervised learning
        always uses ground truth).  At inference, no ground truth
        is available or used.
  \item The \texttt{.detach()} prevents the gradient from flowing
        back through the mask construction, which would create
        a circular dependency.
  \item Analogy: it is similar to importance sampling in RL, where
        the sampling distribution depends on the target but gradients
        are blocked through the weights.
  \item However, this does make the loss non-differentiable with
        respect to the mask shape, so the network cannot learn to
        ``game'' the mask.
\end{itemize}
\textbf{Thesis references:} Sec.~3.7 (loss function definition), Eq.~3.18 (orthogonality loss).
\end{answerbox}

% ── Q19 ─────────────────────────────────────────────────────────────────────
\question{Your asymmetric baseline loss uses $\alpha{=}0.9$ (a 9:1 ratio
penalising over-estimation versus under-estimation).  How sensitive
is the model to this choice?  Did you try $\alpha{=}0.7$ or
$\alpha{=}0.95$?}

\begin{answerbox}
\textbf{Key points:}
\begin{itemize}[nosep]
  \item The intuition for $\alpha{=}0.9$: baseline over-estimation
        absorbs peak energy (leakage), which is harder to correct
        downstream than a small under-estimation.  So we penalise
        over-estimation 9$\times$ more heavily.
  \item Concede: a sensitivity sweep over $\alpha$ was not performed.
        This is a gap.
  \item Argue that the exact value matters less than the
        \emph{qualitative asymmetry}: any $\alpha > 0.5$ biases the
        baseline downward, and the effect saturates.
  \item Very high $\alpha$ (e.g., 0.99) would force the baseline to
        zero everywhere, destroying the baseline estimate.  Very low
        $\alpha$ (e.g., 0.6) provides weak anti-leakage pressure.
  \item A proper sweep (Table~5.4) is listed but not yet run.
\end{itemize}
\textbf{Thesis references:} Sec.~3.7 (asymmetric loss), Table~5.4 (sensitivity ablation, currently empty).
\end{answerbox}

% ── Q20 ─────────────────────────────────────────────────────────────────────
\question{The envelope constraint uses a soft-minimum with temperature
$\tau{=}0.1$ and a window size of 51.  How sensitive are the results
to these choices?  Why 51?}

\begin{answerbox}
\textbf{Key points:}
\begin{itemize}[nosep]
  \item Window size 51 corresponds to roughly the half-width of a
        typical synthetic peak ($\sigma \approx 20$--30 samples), so
        the envelope captures the local minimum in the valley between
        adjacent peaks.
  \item Temperature $\tau{=}0.1$ makes the soft-min approximate a
        hard min: $\text{softmin}_\tau(\mathbf{x}) \to \min(\mathbf{x})$
        as $\tau \to 0$.  At $\tau{=}0.1$, the approximation is tight
        while remaining differentiable.
  \item Sensitivity: smaller windows ($\sim$21) would miss wide peaks;
        larger windows ($\sim$101) would create a baseline that cannot
        follow moderate-frequency drift.
  \item This was not swept systematically.  The values were set once
        based on the known peak widths in the synthetic data.
  \item On real data with unknown peak widths, these would need
        adaptive selection---a limitation.
\end{itemize}
\textbf{Thesis references:} Sec.~3.7 (envelope constraint definition).
\end{answerbox}

% ── Q21 ─────────────────────────────────────────────────────────────────────
\question{You describe the $\ell_1$ and TV penalties as \emph{both} the
cause of baseline leakage \emph{and} necessary for sparsity in
non-peak regions.  This seems like a fundamental contradiction.
How do you resolve it?}

\begin{answerbox}
\textbf{Key points:}
\begin{itemize}[nosep]
  \item This is the \textbf{central tension of the thesis}---do not
        try to resolve it cleanly, because it \emph{is not resolved}.
  \item $\ell_1$ promotes sparsity in $\mathbf{p}$, which is correct
        in regions with no peaks (sparse support).  But in dense-peak
        regions, $\ell_1$ biases peak amplitudes toward zero
        (\emph{soft-thresholding shrinkage bias}, well-known in
        compressed sensing).
  \item The energy removed from peaks has to go somewhere in the
        $\mathbf{y} = \mathbf{b} + \mathbf{p} + \boldsymbol{\varepsilon}$
        decomposition---it leaks into $\hat{\mathbf{b}}$.
  \item The 12-term loss attempts to \emph{mitigate} this with
        anti-leakage terms, but BLI plateaus at $\sim$0.59, confirming
        that loss engineering cannot overcome the architectural bias.
  \item \textbf{The honest conclusion}: $\ell_1$-based unrolling is
        fundamentally limited for dense-peak signals.  Future work
        should explore group sparsity, $\ell_0$-like penalties, or
        hybrid architectures.
\end{itemize}
\textbf{Thesis references:} Sec.~3.7 (loss design rationale), Sec.~5.2.2 (BLI plateau analysis), Sec.~5.5 (leakage discussion).
\end{answerbox}

% ── Q22 ─────────────────────────────────────────────────────────────────────
\question{You use Huber loss with $\delta{=}1.0$ for reconstruction in the
main experiments but mention $\delta{=}0.1$ elsewhere.  Which is it,
and why Huber over MSE in the first place?}

\begin{answerbox}
\textbf{Key points:}
\begin{itemize}[nosep]
  \item Huber loss is piecewise: quadratic for $|e| \le \delta$,
        linear for $|e| > \delta$.  It down-weights large residuals,
        providing robustness to outlier peaks.
  \item $\delta{=}1.0$: the data is normalised to $[0, 1]$, so
        $\delta{=}1.0$ means \emph{almost all residuals fall in the
        quadratic regime}---Huber$_{1.0}$ $\approx$ MSE for
        normalised data.
  \item $\delta{=}0.1$: makes the linear regime much larger, acting
        more like MAE.  This was used in the curriculum's Stage~B
        (Sec.~3.8) to reduce sensitivity to large leakage residuals
        during the early anti-leakage phase.
  \item Acknowledge that the $\delta$ value is not consistently
        documented and should be clarified.
  \item MSE was not used because early experiments showed that
        squared large residuals from leakage regions dominated the
        gradient, destabilising training.
\end{itemize}
\textbf{Thesis references:} Sec.~3.7 (loss terms), Sec.~3.8 (curriculum training), Table~4.2 (training config).
\end{answerbox}


% ============================================================================
%  5. EXPERIMENTAL METHODOLOGY
% ============================================================================
\section{Experimental Methodology}

% ── Q23 ─────────────────────────────────────────────────────────────────────
\question{200 signals total, 40 for testing.  How can you draw
\emph{any} statistically meaningful conclusions from 40 test
signals?  What is your confidence interval on, say, peak
correlation?}

\begin{answerbox}
\textbf{Key points:}
\begin{itemize}[nosep]
  \item Concede: 40 test signals is small.  No confidence intervals
        or statistical tests are reported, which is a limitation.
  \item Mitigating factors: each signal contains 8--25 peaks, so the
        effective sample size for \emph{per-peak} metrics is $\sim$600.
        Per-signal metrics are averaged over $\sim$40 diverse signals.
  \item The synthetic data generation is controlled, so there is no
        domain shift between train and test---only random variation
        in peak parameters.
  \item The small dataset is \emph{by design}: it tests whether
        unrolling's inductive bias can compensate for data scarcity.
        A 100k-signal dataset would not test this hypothesis.
  \item For a proper statistical treatment, report mean $\pm$ std
        and bootstrap confidence intervals.  This should be added
        before the final submission.
\end{itemize}
\textbf{Thesis references:} Sec.~4.1 (data generation), Sec.~5.1 (evaluation metrics), Table~5.1.
\end{answerbox}

% ── Q24 ─────────────────────────────────────────────────────────────────────
\question{Your comparison with classical BEADS uses parameters designed for
unnormalised data, applied to normalised data.  Isn't this a
deliberately unfair comparison that makes BEADS look worse than it
is?}

\begin{answerbox}
\textbf{Key points:}
\begin{itemize}[nosep]
  \item Partially valid criticism.  The default BEADS parameters
        ($\lambda_0{=}1, \lambda_1{=}1, \lambda_2{=}1$) are
        recommended for unnormalised chromatograms with peak heights
        of $10^3$--$10^6$.  On $[0,1]$-normalised data, these
        regularisation strengths are relatively too large.
  \item However: the comparison is intended to show what happens
        when a user applies BEADS ``out of the box'' without
        per-signal tuning.  This is the common use case.
  \item A \emph{fairer} comparison would grid-search BEADS parameters
        on a validation set.  This was not done and should be
        acknowledged as a limitation.
  \item The main contribution is not ``LBEADS-NET beats BEADS'' but
        rather the analysis of leakage and the unrolling methodology.
\end{itemize}
\textbf{Thesis references:} Sec.~5.2 (comparison setup), Table~5.1 (results).
\end{answerbox}

% ── Q25 ─────────────────────────────────────────────────────────────────────
\question{You survey arPLS, airPLS, SNIP, and several deep learning
baselines in Chapter~2, but you never compare against \emph{any}
of them experimentally.  Why not?}

\begin{answerbox}
\textbf{Key points:}
\begin{itemize}[nosep]
  \item Concede: this is a significant omission.
  \item The reason is scope and time: implementing each baseline
        with proper hyperparameter tuning on the synthetic dataset
        would require substantial additional work.
  \item The thesis's primary goal is to unroll \emph{BEADS
        specifically} (given the advisor's authorship of BEADS),
        not to conduct a broad benchmark.
  \item arPLS is the most important missing comparison because it is
        the most widely used method in practice.  This should be
        added before the defense if possible.
  \item Deep learning baselines (1-D U-Net) would require
        architectural design and training, which is a separate study.
  \item Frame as future work.
\end{itemize}
\textbf{Thesis references:} Sec.~2.3 (method survey), Table~2.1 (comparison table), Sec.~6 (future work).
\end{answerbox}

% ── Q26 ─────────────────────────────────────────────────────────────────────
\question{Your synthetic data uses bi-Gaussian peaks (Section~4.1.1), but
your background chapter discusses the EMG (Exponentially Modified
Gaussian) model in detail (Section~2.1.4).  Why the mismatch?  Does
this limit the relevance of your results?}

\begin{answerbox}
\textbf{Key points:}
\begin{itemize}[nosep]
  \item EMG is the standard physical model for chromatographic peaks
        (mass transfer kinetics cause exponential tailing).
        Bi-Gaussian is a simpler empirical model that captures
        asymmetry without the convolution integral.
  \item Bi-Gaussian was chosen for computational simplicity in
        synthetic data generation and because it has an analytic
        closed-form (no numerical convolution needed).
  \item The two models produce qualitatively similar peak shapes
        for moderate asymmetry.  For heavy tailing, EMG produces
        longer tails that bi-Gaussian cannot capture.
  \item This mismatch is not explained in the thesis and should be.
        Add a brief justification in Sec.~4.1.1.
  \item The choice does limit relevance: results on bi-Gaussian
        peaks may not transfer to heavily tailed EMG peaks.
\end{itemize}
\textbf{Thesis references:} Sec.~2.1.4 (EMG model), Sec.~4.1.1 (bi-Gaussian generation).
\end{answerbox}

% ── Q27 ─────────────────────────────────────────────────────────────────────
\question{Your synthetic data generation has many design choices: cluster
probability 0.4, log-uniform amplitude distributions, specific
$\sigma$ ranges.  How sensitive are your results to these choices?
Have you tested on a different distribution?}

\begin{answerbox}
\textbf{Key points:}
\begin{itemize}[nosep]
  \item These parameters control the ``difficulty distribution'' of
        the training set: cluster probability determines how often
        dense-peak regions occur (which is where leakage is worst).
  \item No sensitivity analysis on data generation parameters was
        performed.  This is a limitation.
  \item Higher cluster probability would create more dense-peak
        signals, likely increasing average BLI.  Lower cluster
        probability would create easier signals with better metrics.
  \item The log-uniform amplitude distribution ensures coverage of
        both small and large peaks, which is more realistic than
        uniform sampling.
  \item A robust model should be tested on out-of-distribution
        data (different peak width ranges, different noise levels).
        This was not done.
\end{itemize}
\textbf{Thesis references:} Sec.~4.1.1 (data generation parameters), Table~4.1 (generation config).
\end{answerbox}


% ============================================================================
%  6. RESULTS INTERPRETATION
% ============================================================================
\section{Results Interpretation}

% ── Q28 ─────────────────────────────────────────────────────────────────────
\question{BLI is approximately 0.59 across \emph{all} configurations---from
the simplest 4-term loss to the full 12-term v7 loss.  Doesn't this
mean your entire loss engineering effort was basically useless for
addressing leakage?}

\begin{answerbox}
\textbf{Key points:}
\begin{itemize}[nosep]
  \item \textbf{Yes, essentially.}  This is a key negative result
        and should be stated directly.
  \item The BLI plateau at $\sim$0.59 across all configurations is
        strong evidence that leakage is an \textbf{architectural}
        property of $\ell_1$-based unrolling, not a loss-function
        problem.
  \item This is a \emph{valuable finding}: it tells future
        researchers not to invest in loss engineering for this
        architecture.  The fix must be structural (different penalty,
        different architecture).
  \item The loss terms \emph{do} help other metrics (reconstruction
        error, peak shape fidelity), just not leakage.
  \item Frame this as: ``the thesis \emph{discovered} the leakage
        floor and \emph{proved} that loss engineering cannot break
        through it.''
\end{itemize}
\textbf{Thesis references:} Table~5.1 (BLI column), Sec.~5.2.2 (BLI analysis), Sec.~5.5 (leakage discussion).
\end{answerbox}

% ── Q29 ─────────────────────────────────────────────────────────────────────
\question{Area error is 49--58\% across all LBEADS-NET variants.  In
pharmaceutical QC, 1\% accuracy is the requirement.  50\% error
is catastrophic.  How can you call this a contribution?}

\begin{answerbox}
\textbf{Key points:}
\begin{itemize}[nosep]
  \item Do not try to defend 50\% area error as acceptable---it is
        not, and any committee member from chromatography would dismiss
        this immediately.
  \item Reframe: the contribution is \emph{not} a production-ready
        tool.  It is a \textbf{methodological study} that:
        (a)~demonstrates algorithm unrolling for BEADS,
        (b)~identifies and characterises the $\ell_1$ leakage limit,
        (c)~documents the failure of loss engineering to fix it.
  \item The area error is dominated by leakage in dense-peak regions.
        For isolated peaks, the per-peak area error is much lower
        (this should be reported separately if data is available).
  \item The hybrid pipeline (Sec.~3.9) improves area error by
        applying classical BEADS post-processing, but the fundamental
        limit remains.
  \item Honest conclusion: LBEADS-NET in its current form is not
        suitable for quantitative chromatography.
\end{itemize}
\textbf{Thesis references:} Table~5.1 (area error), Sec.~5.2 (quantitative results), Sec.~5.5.
\end{answerbox}

% ── Q30 ─────────────────────────────────────────────────────────────────────
\question{Experiment~1.5 (the full v7 12-term loss) performs \emph{worse}
than Experiment~1.3 (a simpler configuration).  Your own results
show that simpler is better.  Why do you propose the complex
version at all?}

\begin{answerbox}
\textbf{Key points:}
\begin{itemize}[nosep]
  \item This is a genuinely important finding.  The 12-term loss
        creates \textbf{competing gradient directions}: the
        reconstruction loss pulls parameters one way while the
        anti-leakage losses pull them another way, resulting in
        a compromise that is worse on all metrics.
  \item This is an example of \emph{over-regularisation}---adding
        too many constraints to a 100-parameter model.
  \item The simpler Experiment~1.3 succeeds because fewer loss terms
        give the optimiser a cleaner gradient signal.
  \item The thesis proposes v7 as a \emph{research contribution}
        (novel loss formulations) even though v5/Exp.~1.3 is
        empirically better.  The intellectual value is in
        understanding \emph{why} the complex version fails.
  \item Recommendation: the thesis should more prominently highlight
        Exp.~1.3 as the best-performing configuration.
\end{itemize}
\textbf{Thesis references:} Table~5.1 (Exp.~1.3 vs 1.5), Sec.~5.2.1 (progressive results).
\end{answerbox}

% ── Q31 ─────────────────────────────────────────────────────────────────────
\question{Peak correlation maxes at 0.793.  What does this mean physically?
That 79\% of the peak shape is recovered?  What is an acceptable
value in practice?}

\begin{answerbox}
\textbf{Key points:}
\begin{itemize}[nosep]
  \item Correlation measures \emph{shape similarity} (cosine of the
        angle between the estimated and true peak vectors), not
        amplitude accuracy.  A correlation of 0.793 means the
        estimated peaks preserve $\sim$79\% of the shape structure.
  \item The remaining $\sim$21\% discrepancy comes from
        (a)~amplitude shrinkage due to $\ell_1$ bias and
        (b)~baseline leakage distorting the peak shape, especially
        at the base.
  \item In practice, $r > 0.95$ is expected for quantitative
        chromatography.  $0.793$ would not be acceptable.
  \item For qualitative analysis (peak detection, retention time
        identification), $0.793$ may be sufficient since the peak
        \emph{positions} are typically preserved even when amplitudes
        are distorted.
  \item The gap from 0.793 to 1.0 is another manifestation of the
        $\ell_1$ leakage problem.
\end{itemize}
\textbf{Thesis references:} Sec.~5.1 (metric definitions), Table~5.1 (peak correlation).
\end{answerbox}

% ── Q32 ─────────────────────────────────────────────────────────────────────
\question{The learned parameter analysis (Section~5.2.3) uses a 6-layer
model with $N{=}1024$, but the main results use a 20-layer model
with $N{=}4096$.  Are these comparable?  Why not analyse the
parameters of the actual model you evaluate?}

\begin{answerbox}
\textbf{Key points:}
\begin{itemize}[nosep]
  \item They are \textbf{not directly comparable}, and this should
        be acknowledged.
  \item The 6-layer/$N{=}1024$ model was used for the parameter
        analysis because (a)~it is easier to visualise 6 layers
        and (b)~it was an earlier experiment before the full
        20-layer model was trained.
  \item The qualitative pattern (early layers learn coarse separation,
        later layers learn refinement) is expected to hold across
        depths, but the specific parameter values would differ.
  \item Ideally, the parameter trajectories of the 20-layer model
        should be plotted.  This is a gap.
  \item The analysis still provides insight into the
        \emph{specialisation phenomenon}, which is the main finding
        of that section, even if the numerical values are
        model-specific.
\end{itemize}
\textbf{Thesis references:} Sec.~5.2.3 (parameter analysis), Sec.~4.2 (model configurations).
\end{answerbox}


% ============================================================================
%  7. BASELINE LEAKAGE DEEP DIVE
% ============================================================================
\section{Baseline Leakage Deep Dive}

% ── Q33 ─────────────────────────────────────────────────────────────────────
\question{You claim leakage is a ``fundamental limit of $\ell_1$-based
formulations.''  But classical BEADS with properly tuned parameters
on unnormalised data presumably has lower leakage.  Isn't the
leakage just a parameter-tuning problem that you failed to solve?}

\begin{answerbox}
\textbf{Key points:}
\begin{itemize}[nosep]
  \item The leakage mechanism is \emph{not} a tuning artefact:
        it stems from the soft-thresholding operator
        $\text{soft}(x, \lambda) = \text{sign}(x)(|x| - \lambda)_+$,
        which \emph{always} shrinks non-zero entries by exactly
        $\lambda$.  This bias exists for any $\lambda > 0$.
  \item Classical BEADS on unnormalised data has the \emph{same}
        mechanism, but the effect may be less visible because
        (a)~peak amplitudes are $10^3$--$10^6$ so the relative
        shrinkage $\lambda / A$ is small, and (b)~users tune
        $\lambda$ to minimise visible artefacts.
  \item On normalised $[0,1]$ data, the shrinkage $\lambda / A$
        is relatively larger, making leakage more prominent.
  \item The deeper point: in \emph{dense-peak regions}, the
        sparsity assumption fails entirely (peaks overlap, the
        signal is not sparse), so $\ell_1$ is the wrong prior.
  \item The leakage is thus both a tuning issue (partially)
        \emph{and} a structural issue (fundamentally, in dense regions).
\end{itemize}
\textbf{Thesis references:} Sec.~5.5 (leakage analysis), Sec.~3.5 (proximal operator), Eq.~3.11.
\end{answerbox}

% ── Q34 ─────────────────────────────────────────────────────────────────────
\question{You mention group sparsity as a potential solution to leakage.
Can you explain concretely how it would help, and why you did not
pursue it?}

\begin{answerbox}
\textbf{Key points:}
\begin{itemize}[nosep]
  \item Group sparsity (e.g., group LASSO or $\ell_{2,1}$ penalty)
        penalises \emph{groups} of adjacent coefficients jointly:
        $\sum_g \|\mathbf{p}_g\|_2$.  Either the entire group is
        zero or non-zero---no element-wise shrinkage within active
        groups.
  \item This directly addresses the leakage mechanism: within an
        active peak group, amplitudes are \emph{not} shrunk toward
        zero, preserving peak shape and area.
  \item Implementation challenge: defining the groups.  Peak
        boundaries are unknown a priori, so the grouping must be
        adaptive or overlapping (e.g., sliding-window group LASSO).
  \item The proximal operator for group $\ell_{2,1}$ is a
        \emph{block} soft-thresholding, which is straightforward
        to implement in PyTorch.
  \item Not pursued due to scope and time constraints.  This is
        the highest-priority future work.
\end{itemize}
\textbf{Thesis references:} Sec.~5.5 (leakage discussion), Sec.~6 (future work).
\end{answerbox}

% ── Q35 ─────────────────────────────────────────────────────────────────────
\question{Gharbi et al.'s U-HQ method uses a hybrid penalty that avoids
$\ell_1$ shrinkage bias.  You cite it in the background.  Why not
adopt their approach instead of fighting the $\ell_1$ bias with
loss terms?}

\begin{answerbox}
\textbf{Key points:}
\begin{itemize}[nosep]
  \item U-HQ (Gharbi et al.) uses a non-convex penalty that
        approximates $\ell_0$: it does not shrink large coefficients,
        only eliminates small ones.  This would indeed avoid leakage.
  \item The challenge: non-convex penalties lose the convergence
        guarantees of MM/ISTA.  The proximal operator may not be
        unique, and unrolling a non-convergent algorithm is
        theoretically questionable.
  \item Adopting U-HQ would require redesigning the entire unrolling
        framework: different cost function, different proximal operator,
        different convergence analysis.
  \item This is the right direction for future work, but it is a
        separate thesis.
  \item The contribution here is establishing \emph{why} the $\ell_1$
        version fails, which motivates the move to U-HQ-style penalties.
\end{itemize}
\textbf{Thesis references:} Sec.~2.3 (U-HQ discussion), Sec.~5.5 (leakage analysis), Sec.~6.
\end{answerbox}

% ── Q36 ─────────────────────────────────────────────────────────────────────
\question{If baseline leakage is truly a fundamental limitation of the
$\ell_1$ formulation, then why build on BEADS at all?  Isn't
choosing a flawed base algorithm the core problem?}

\begin{answerbox}
\textbf{Key points:}
\begin{itemize}[nosep]
  \item BEADS was chosen because (a)~it is the advisor's algorithm,
        making this a natural extension of the group's research, and
        (b)~it is one of the most widely used baseline estimation
        methods with a clear optimisation structure suitable for
        unrolling.
  \item BEADS is not ``flawed''---it is a well-designed convex
        programme that works well for sparse-peak chromatograms.
        The limitation arises specifically in the dense-peak regime.
  \item The thesis demonstrates \emph{exactly where} BEADS breaks
        down when unrolled.  This is a contribution: understanding
        the limits of an existing method rigorously.
  \item The negative result is arguably more valuable than a modest
        positive improvement: it redirects future research toward
        structural solutions (group sparsity, non-convex penalties).
  \item ``All models are wrong; some are useful.''  BEADS is useful;
        we now know precisely where it stops being useful.
\end{itemize}
\textbf{Thesis references:} Sec.~1.2 (thesis scope), Sec.~2.3.3 (BEADS), Sec.~5.5 (limitations).
\end{answerbox}


% ============================================================================
%  8. PRACTICAL / "SO WHAT" QUESTIONS
% ============================================================================
\section{Practical / ``So What'' Questions}

% ── Q37 ─────────────────────────────────────────────────────────────────────
\question{Has LBEADS-NET been validated on even \emph{one} real
chromatographic dataset with known ground truth?}

\begin{answerbox}
\textbf{Key points:}
\begin{itemize}[nosep]
  \item \textbf{No.}  All experiments use synthetic data.
  \item The fundamental problem: real chromatograms do not have
        ground-truth decompositions ($\mathbf{b}$, $\mathbf{p}$,
        $\boldsymbol{\varepsilon}$ are not separately measurable).
  \item Semi-synthetic validation is possible: take a known peak
        standard (e.g., caffeine), subtract a fitted baseline
        manually, and use the result as ``ground truth.''  This
        was not attempted.
  \item Public chromatographic datasets with ground truth are
        essentially non-existent.  This is a known challenge in
        the field.
  \item Real-data validation is the most important next step (Sec.~6).
\end{itemize}
\textbf{Thesis references:} Sec.~4.1 (synthetic data only), Sec.~6 (future work).
\end{answerbox}

% ── Q38 ─────────────────────────────────────────────────────────────────────
\question{What is the inference time of LBEADS-NET?  If each layer requires
a dense $N \times N$ matrix--vector multiply, it might be slower
than classical BEADS, which is $O(N)$ per iteration with banded
solvers.}

\begin{answerbox}
\textbf{Key points:}
\begin{itemize}[nosep]
  \item Inference time is \textbf{not reported} in the thesis---this
        is a gap.
  \item Each of the $K{=}20$ layers performs $O(N^2)$ operations
        due to dense matrix multiplication.  For $N{=}4096$, this
        is $\sim$20 $\times$ $16\text{M}$ = $320\text{M}$ flops.
  \item Classical BEADS with banded solvers is $O(N \cdot b \cdot K)$
        where $b$ is bandwidth ($\sim$5--10), so $\sim$20 $\times$
        $40\text{K}$ = $800\text{K}$ flops---\textbf{400$\times$ fewer}.
  \item On a GPU, the dense matmuls are parallelised, so wall-clock
        time may be comparable.  But on CPU (the typical analytical
        lab environment), LBEADS-NET is significantly slower.
  \item With banded implementation, LBEADS-NET would also be $O(N)$
        per layer, eliminating the gap.
\end{itemize}
\textbf{Thesis references:} Sec.~3.4 (implementation), Sec.~6 (future work on banded ops).
\end{answerbox}

% ── Q39 ─────────────────────────────────────────────────────────────────────
\question{How would a chromatographer actually use LBEADS-NET?  They do not
have ground-truth baselines for training.  Does the model
generalise to unseen instrument configurations?}

\begin{answerbox}
\textbf{Key points:}
\begin{itemize}[nosep]
  \item Usage scenario: train on synthetic data that covers the
        expected range of peak shapes, baseline curvatures, and
        noise levels; deploy on real data.
  \item This is the standard paradigm for model-based deep learning
        in domains without ground truth (e.g., MRI reconstruction).
  \item Generalisation to unseen instruments depends on whether
        the synthetic data distribution covers the real data
        distribution.  This is not tested.
  \item Domain adaptation (fine-tuning on a small set of manually
        corrected real chromatograms) is a natural extension.
  \item Practically, the user would run LBEADS-NET as a preprocessing
        step in their chromatography data system (CDS) software,
        replacing the built-in baseline correction module.
\end{itemize}
\textbf{Thesis references:} Sec.~4.1 (synthetic data rationale), Sec.~6 (deployment discussion).
\end{answerbox}

% ── Q40 ─────────────────────────────────────────────────────────────────────
\question{Your hybrid pipeline adds classical BEADS as a post-processing
step.  If you still need classical BEADS at the end, what did
unrolling actually buy you?}

\begin{answerbox}
\textbf{Key points:}
\begin{itemize}[nosep]
  \item The hybrid pipeline uses LBEADS-NET to provide an
        \emph{initial decomposition} that is closer to the true
        solution than the default initialisation ($\mathbf{b}_0 = 0$).
        Classical BEADS then refines from this warm start.
  \item The benefit: BEADS converges in fewer iterations from a
        good initialisation, and the learned parameters provide
        signal-adaptive regularisation weights that BEADS alone
        cannot.
  \item Concede: if classical BEADS with proper tuning achieves
        similar results, the unrolling step adds complexity without
        clear benefit.
  \item The honest framing: the hybrid pipeline is a
        \emph{fallback strategy} when LBEADS-NET alone is insufficient,
        not the primary contribution.
  \item The primary contribution is the understanding of leakage,
        not the pipeline itself.
\end{itemize}
\textbf{Thesis references:} Sec.~3.9 (hybrid pipeline), Sec.~5.2 (hybrid results).
\end{answerbox}


% ============================================================================
%  9. WRITING & COMPLETENESS
% ============================================================================
\section{Writing \& Completeness}

% ── Q41 ─────────────────────────────────────────────────────────────────────
\question{Your abstract is Lorem Ipsum.  Your conclusions chapter is Lorem
Ipsum.  Your ablation tables have all dashes.  Several sections are
marked as placeholders.  Is this thesis actually complete?}

\begin{answerbox}
\textbf{Key points:}
\begin{itemize}[nosep]
  \item \textbf{Do not get defensive.}  Acknowledge that these
        sections are incomplete and will be finalised before the
        official submission deadline.
  \item Emphasise that the \emph{technical content} (Chapters 2--5)
        is complete and the defence focuses on the technical
        contribution.
  \item Provide a timeline: abstract and conclusions will be written
        after the defence, incorporating committee feedback.
  \item For the ablation tables: state which experiments are planned
        and the expected timeline.
  \item This is a common situation for defences that occur before
        the final manuscript submission.
\end{itemize}
\textbf{Thesis references:} Abstract, Chapter~6 (conclusions), Tables~5.3--5.4.
\end{answerbox}

% ── Q42 ─────────────────────────────────────────────────────────────────────
\question{Several cross-references in your document show ``Chapter~??''
(pages 43, 56, 68).  Have you compiled and checked the final PDF?}

\begin{answerbox}
\textbf{Key points:}
\begin{itemize}[nosep]
  \item Acknowledge the broken references.  These are caused by
        missing \verb|\label{}| tags or mismatched \verb|\ref{}| keys.
  \item State that a full compilation pass with \verb|latexmk|
        (multiple runs to resolve references) will be done before
        submission.
  \item This is a mechanical fix, not a content issue.
  \item Have a plan: grep for ``??'' in the compiled PDF, trace each
        to the source \verb|\ref{}| command, and fix.
\end{itemize}
\textbf{Thesis references:} Throughout document (cross-references).
\end{answerbox}

% ── Q43 ─────────────────────────────────────────────────────────────────────
\question{You say ``twelve-term composite loss'' in Section~3.7 and
Table~3.3, but your Introduction says ``11-term.''  Which is
correct?}

\begin{answerbox}
\textbf{Key points:}
\begin{itemize}[nosep]
  \item The correct count is \textbf{twelve}.  Section~3.7 and
        Table~3.3 enumerate all twelve terms explicitly.
  \item The Introduction's ``11-term'' is a typo from an earlier
        version of the loss function before the twelfth term
        (likely the envelope constraint) was added.
  \item This will be corrected in the Introduction before
        submission.
  \item Having the specific section and table to point to
        (Sec.~3.7, Table~3.3) provides a definitive answer.
\end{itemize}
\textbf{Thesis references:} Sec.~1.2 (Introduction, ``11-term''), Sec.~3.7 (``twelve-term''), Table~3.3.
\end{answerbox}


% ============================================================================
%  10. FUTURE WORK & BIG PICTURE
% ============================================================================
\section{Future Work \& Big Picture}

% ── Q44 ─────────────────────────────────────────────────────────────────────
\question{If you had 6 more months, what would you do first and why?}

\begin{answerbox}
\textbf{Key points (prioritised):}
\begin{enumerate}[nosep]
  \item \textbf{Group sparsity penalty} (2 months): replace $\ell_1$
        with $\ell_{2,1}$ or overlapping group LASSO.  This directly
        addresses the leakage ceiling and is the highest-impact change.
  \item \textbf{Banded matrix implementation} (1 month): exploit the
        sparsity structure of $\mathbf{H}$ and $\mathbf{D}$ to reduce
        memory from $O(N^2)$ to $O(N)$ and enable $N{=}65536$+ signals.
  \item \textbf{Real data validation} (1 month): obtain chromatographic
        data from a collaborating lab, create semi-synthetic ground
        truth, and evaluate.
  \item \textbf{arPLS/airPLS baseline comparison} (2 weeks): implement
        and tune arPLS on the existing synthetic dataset.
  \item \textbf{Learnable $f_c$} (2 weeks): parameterise the highpass
        filter as a learnable FIR filter initialised from the
        Butterworth design.
  \item \textbf{Complete ablations} (2 weeks): fill Tables~5.3 and
        5.4.
\end{enumerate}
\textbf{Thesis references:} Sec.~6 (future work), Sec.~5.5 (leakage analysis).
\end{answerbox}

% ── Q45 ─────────────────────────────────────────────────────────────────────
\question{Could the LBEADS-NET approach generalise to other signal
decomposition problems---Raman spectroscopy, NMR, mass
spectrometry?}

\begin{answerbox}
\textbf{Key points:}
\begin{itemize}[nosep]
  \item \textbf{Yes, in principle.}  Any problem where the signal
        decomposes as $\mathbf{y} = \mathbf{b} + \mathbf{p} +
        \boldsymbol{\varepsilon}$ with a smooth baseline and sparse
        peaks is a candidate.
  \item \textbf{Raman spectroscopy}: fluorescence background is
        smooth and broadband, peaks are narrow Lorentzian---very
        similar structure.  BEADS-like methods are already used
        (Zhang et al., 2010).
  \item \textbf{NMR}: baseline distortion from phase errors and
        receiver artefacts.  The peak model is Lorentzian with
        known analytic form.
  \item \textbf{Mass spectrometry}: chemical noise creates a
        baseline, but the signal is often discrete (centroid data),
        which changes the problem structure.
  \item The unrolling framework is modular: replace the peak prior
        (Lorentzian instead of bi-Gaussian) and the filter design
        (adjust $f_c$ for the spectral content), and retrain.
  \item The $\ell_1$ leakage limitation would apply equally in
        all these domains---the proposed group sparsity fix would
        also be needed.
\end{itemize}
\textbf{Thesis references:} Sec.~2.1 (signal model), Sec.~6 (future work).
\end{answerbox}

% ── Q46 ─────────────────────────────────────────────────────────────────────
\question{What would it take to make $f_c$ learnable?  Sketch the approach
at the whiteboard.}

\begin{answerbox}
\textbf{Key points (sketch for the whiteboard):}
\begin{enumerate}[nosep]
  \item \textbf{Current}: $f_c$ enters through
        \texttt{scipy.signal.butter(d, f\_c)} $\to$ IIR coefficients
        $\to$ dense matrix $\mathbf{H}$.  This pipeline is not
        differentiable.
  \item \textbf{Approach 1---Differentiable IIR:} Place Butterworth
        poles at $p_k = e^{j(\pi/2 \pm \Delta\theta_k)}$ where
        $\Delta\theta_k = g(f_c)$.  Construct the transfer function
        $H(z) = \prod_k (1 - p_k z^{-1})$ in PyTorch.  Apply via
        \texttt{torch.linalg.solve\_triangular} (causal IIR).
        Gradients flow through $f_c \to p_k \to H(z) \to
        \mathbf{H}\mathbf{y}$.
  \item \textbf{Approach 2---Learnable FIR:} Design an FIR filter
        of order $M$ initialised from the Butterworth impulse
        response truncated at $M$ taps.  Make the tap weights
        $h_0, \ldots, h_M$ learnable.  Apply via
        \texttt{torch.nn.functional.conv1d}.  This is simpler but
        loses the sharp IIR rolloff.
  \item \textbf{Approach 3---Frequency-domain:} Multiply
        $\mathbf{Y}(f)$ by a learnable frequency response
        $|\hat{H}(f)|$ parameterised by $f_c$ and filter order.
        Invert via IFFT.  Fully differentiable but introduces
        circular convolution artefacts.
  \item \textbf{Recommended}: Approach~2 (FIR) for simplicity and
        robustness, with $M \sim 64$ taps.
\end{enumerate}
\textbf{Thesis references:} Sec.~3.4 (current filter construction), Sec.~6 (future work).
\end{answerbox}


% ============================================================================
%  11. KILLER QUESTIONS -- THE MEANEST ONES
% ============================================================================
\newpage
\section{Killer Questions --- The Ones That Hurt}

\begin{tcolorbox}[
  enhanced,
  colback=WarnBG,
  colframe=WarnFrame,
  fonttitle=\bfseries\sffamily,
  title={Warning},
  coltitle=white,
  attach boxed title to top left={yshift=-2mm, xshift=4mm},
  boxed title style={colback=WarnFrame, colframe=WarnFrame, sharp corners, size=small},
  left=4pt, right=4pt, top=4pt, bottom=4pt
]
These are the questions designed to make you uncomfortable.  They target
circular reasoning, contribution validity, mathematical rigour, and the
``so what'' factor.  If you can handle these, you can handle anything.
\end{tcolorbox}

\vspace{0.5em}

% ── Q47 ─────────────────────────────────────────────────────────────────────
\question{Strip away the loss engineering that your own results show does
not work.  What remains?  You unrolled BEADS into ISTA layers with
scalar parameters---that is just LISTA applied to a different cost
function.  Where, exactly, is the novelty?}

\begin{answerbox}
\textbf{Key points:}
\begin{itemize}[nosep]
  \item Acknowledge the similarity to LISTA, but clarify the differences:
        (a)~LISTA learns full matrices $\mathbf{W}_e, \mathbf{S}$;
        LBEADS-NET fixes the operators and learns only scalars.
        (b)~LISTA targets sparse coding; LBEADS-NET targets a
        \emph{three-component} signal decomposition with frequency-domain
        separation (baseline $+$ peaks $+$ noise).
  \item The novelty is application-specific:
        \emph{no prior work has unrolled the BEADS algorithm}.
        The ISTA reformulation of the MM update (replacing CG inner
        solve with proximal gradient) is a design contribution.
  \item The identification and systematic characterisation of
        baseline leakage as a fundamental $\ell_1$ limitation in
        this architecture is novel.
  \item The 12-term anti-leakage loss design (even though it fails
        to break the BLI floor) and the curriculum strategy are
        original engineering contributions.
  \item Concede: the theoretical novelty is modest.  This is an
        \emph{applied} thesis.
\end{itemize}
\textbf{Thesis references:} Sec.~2.3.2 (LISTA), Sec.~3.2--3.3 (ISTA reformulation), Sec.~1.2 (contributions).
\end{answerbox}

% ── Q48 ─────────────────────────────────────────────────────────────────────
\question{You train on synthetic data and test on synthetic data generated
by the \emph{same} code with the \emph{same} distributions.  This is
a closed loop.  You have proven nothing about real chromatography.
How do you respond?}

\begin{answerbox}
\textbf{Key points:}
\begin{itemize}[nosep]
  \item This is correct, and it is a fundamental limitation.
  \item The synthetic data was calibrated against 8 real chromatograms
        (Sec.~4.1)---peak heights 10--2600+, widths FWHM 3--370+,
        29--58 peaks per signal---but the test set is still synthetic.
  \item The thesis contribution is \emph{algorithmic}: demonstrating
        unrolling and characterising leakage.  The synthetic evaluation
        framework controls for confounding variables that real data
        would introduce.
  \item All supervised deep learning for baseline correction
        (Kensert et al., Chen et al., Han et al.) uses synthetic
        training data---this is standard practice because ground truth
        does not exist for real chromatograms.
  \item The honest response: real-data validation is the single most
        important missing piece and the top priority for future work.
\end{itemize}
\textbf{Thesis references:} Sec.~4.1 (data generation), Sec.~5.6 (real chromatogram discussion).
\end{answerbox}

% ── Q49 ─────────────────────────────────────────────────────────────────────
\question{Your standard deviations are enormous.  Peak MSE is $9.30
\pm 4.7 \times 10^{-3}$.  That means some signals have MSE of
$4.6 \times 10^{-3}$ and others have $14.0 \times 10^{-3}$---a
$3\times$ range.  Is your model actually learning anything consistent,
or is it just averaging over easy and hard cases?}

\begin{answerbox}
\textbf{Key points:}
\begin{itemize}[nosep]
  \item The high variance is expected: the test set deliberately
        includes signals with very different peak densities (20--55
        peaks), baseline shapes (spline, exponential, mixed), and
        noise levels ($\sigma_w \in [2, 6]$).
  \item A signal with 20 well-separated peaks and low noise is
        genuinely \emph{easier} than one with 55 overlapping peaks
        and high noise.  The variance reflects problem difficulty,
        not model inconsistency.
  \item To verify: stratify results by peak density (sparse vs
        dense) and report separately.  This was not done but should be.
  \item The \emph{relative} improvement over classical BEADS is
        consistent across signals (the std of the improvement would
        be more informative than the std of the raw metric).
  \item Concede: reporting only mean $\pm$ std is insufficient.
        Median, IQR, and per-signal scatter plots would be more
        informative.
\end{itemize}
\textbf{Thesis references:} Table~5.1 (standard deviations), Sec.~4.1 (data diversity).
\end{answerbox}

% ── Q50 ─────────────────────────────────────────────────────────────────────
\question{You say LBEADS-NET has ``only 100 learnable parameters.''  But the
dense matrices $\mathbf{A}$, $\mathbf{B}$, and $\mathbf{L}$ are
each $N \times N = 4096^2 \approx 16.7$ million entries.  These are
frozen, not absent.  The model capacity is dominated by these
fixed operators.  Isn't ``100 parameters'' misleading?}

\begin{answerbox}
\textbf{Key points:}
\begin{itemize}[nosep]
  \item This is a valid nuance.  The \emph{learnable} parameter
        count is 100 scalars.  The \emph{total model capacity}
        includes the fixed filter matrices, which encode substantial
        prior knowledge.
  \item However, the fixed matrices are determined entirely by two
        hyperparameters ($f_c$ and $d$) and the signal length $N$.
        They are not ``learned from data''---they are analytically
        derived from the Butterworth filter design.
  \item The analogy: a CNN with a fixed Gabor filter bank and 100
        learnable classification weights also has ``100 learnable
        parameters'' even though the Gabor filters encode structure.
  \item The ``100 parameters'' claim is about \emph{what backpropagation
        can adjust}, not about the total information content of the
        model.  This distinction should be made clearer in the thesis.
  \item The correct framing: ``100 learnable parameters, with
        fixed structural operators derived from the BEADS
        formulation.''
\end{itemize}
\textbf{Thesis references:} Sec.~3.4 (architecture), Table~3.1 (parameters per layer).
\end{answerbox}

% ── Q51 ─────────────────────────────────────────────────────────────────────
\question{Why not just use Bayesian optimisation to tune the 6 classical
BEADS parameters on a validation set?  That would be simpler than
unrolling, fully interpretable, and might work just as well.}

\begin{answerbox}
\textbf{Key points:}
\begin{itemize}[nosep]
  \item Bayesian optimisation (BO) over $\{f_c, d, r, \lambda_0,
        \lambda_1, \lambda_2\}$ on a validation set is a perfectly
        reasonable baseline.  It was not compared against, and it
        should have been.
  \item Key difference: BO finds a \emph{single global} parameter
        set optimised for the validation distribution.  LBEADS-NET
        learns \emph{per-layer} parameters that specialise across
        depth (Sec.~5.5)---this is fundamentally richer.
  \item BO-tuned BEADS would still suffer from the same $\ell_1$
        leakage in dense-peak regions (the leakage is structural,
        not a tuning issue).  But BO would likely achieve better
        results than the ``unfair'' default-parameter comparison.
  \item Concede: BO-tuned BEADS is the missing comparison that would
        isolate the benefit of \emph{unrolling} from the benefit of
        \emph{better parameters}.
  \item This is arguably the most important missing experiment.
\end{itemize}
\textbf{Thesis references:} Sec.~4.5 (comparison methods), Sec.~5.1 (results).
\end{answerbox}

% ── Q52 ─────────────────────────────────────────────────────────────────────
\question{Your three-stage curriculum training is just a manually designed
training schedule.  You replaced BEADS's 6 hand-tuned parameters
with 12 hand-tuned loss weights \emph{plus} a hand-designed
3-stage curriculum with hand-picked epoch budgets.  How is this
reducing human design effort?}

\begin{answerbox}
\textbf{Key points:}
\begin{itemize}[nosep]
  \item This is a fair and damaging critique.  Acknowledge it.
  \item The key distinction: the manual design effort is a
        \emph{one-time cost} during model development.  Once trained,
        the model requires \emph{zero} user input at inference time.
        Classical BEADS requires per-signal tuning at every use.
  \item The curriculum and loss weights are meta-level design
        decisions analogous to choosing a network architecture in
        deep learning---they are set by the model developer, not
        the end user.
  \item Concede: the total design effort for LBEADS-NET (12 loss
        weights, 3 curriculum stages, epoch budgets, etc.) is
        \emph{greater} than tuning BEADS for a single signal.  The
        payoff is generalisation across signals.
  \item Automated hyperparameter optimisation (e.g., learning the
        loss weights via meta-learning) would strengthen this claim
        and is a natural extension.
\end{itemize}
\textbf{Thesis references:} Sec.~3.6 (curriculum), Table~3.2 (stage configuration), Sec.~3.7 (loss weights).
\end{answerbox}

% ── Q53 ─────────────────────────────────────────────────────────────────────
\question{Gharbi et al.\ already showed in 2024 that $\ell_1$-based
unrolled networks exhibit shrinkage bias in chromatographic
signal recovery.  You cite their work.  Your ``fundamental insight''
about baseline leakage is a confirmation of their finding, not a
new discovery.  What do you actually add beyond their result?}

\begin{answerbox}
\textbf{Key points:}
\begin{itemize}[nosep]
  \item Gharbi et al.\ (2024) observed that $\ell_1$-based unrolled
        methods (U-PD, U-ISTA) produce biased peak intensities
        compared to their U-HQ hybrid penalty.  This is indeed the
        same underlying phenomenon.
  \item What LBEADS-NET adds:
        (a)~\textbf{Quantification}: the BLI and PBER metrics provide
        the first \emph{numerical characterisation} of the leakage
        magnitude (BLI $\approx$ 0.59).
        (b)~\textbf{Systematic mitigation attempt}: 12 anti-leakage
        loss terms, progressive experiments, and the proof that loss
        engineering cannot break the floor.
        (c)~\textbf{Different algorithm}: BEADS (MM-based) vs Gharbi's
        PD/ISTA/HQ formulations.
        (d)~\textbf{Curriculum training}: novel for this domain.
  \item Concede: the conceptual finding is not new.  The
        contribution is the \emph{depth of analysis} and the specific
        characterisation for the BEADS architecture.
  \item Frame as: ``We independently confirmed and quantified a
        phenomenon that Gharbi et al.\ identified qualitatively.''
\end{itemize}
\textbf{Thesis references:} Sec.~2.3.5 (Gharbi et al.), Sec.~5.1.2 (BLI analysis), Sec.~1.3 (contributions).
\end{answerbox}

% ── Q54 ─────────────────────────────────────────────────────────────────────
\question{At the end of the day: your best model has 50\% area error,
0.79 peak correlation, and 0.59 BLI.  A practising chromatographer
would never deploy this.  Your loss engineering does not fix
leakage.  Your ablation tables are empty.  What exactly is the
contribution to knowledge?}

\begin{answerbox}
\textbf{Key points (this is THE critical question---have this answer cold):}
\begin{itemize}[nosep]
  \item \textbf{Contribution 1}: First unrolling of the BEADS
        algorithm.  Demonstrated that BEADS can be reformulated as
        ISTA layers, trained end-to-end, and applied without
        per-signal parameter tuning.
  \item \textbf{Contribution 2}: Identification and quantification
        of the $\ell_1$ leakage floor (BLI $\approx$ 0.59) as a
        \emph{fundamental architectural limit} that loss engineering
        cannot overcome.  This is a negative result with clear
        implications for the field.
  \item \textbf{Contribution 3}: The failure mode catalogue---CG
        gradient vanishing, non-learnable $f_c$, train/inference
        mismatch, softplus artefacts---provides practical guidance
        for anyone attempting algorithm unrolling.
  \item \textbf{Contribution 4}: Emergent parameter specialisation
        across layers, demonstrating that unrolled networks learn
        non-uniform iteration schedules that classical algorithms
        cannot express.
  \item \textbf{The framing}: this is a \emph{methodological
        exploration} and \emph{negative result study}, not a
        deployment-ready tool.  The value is in \emph{what we learned}
        about unrolling sparsity-based algorithms.
\end{itemize}
\textbf{Thesis references:} Sec.~1.3 (contributions), Sec.~5.7 (summary of findings).
\end{answerbox}

% ── Q55 ─────────────────────────────────────────────────────────────────────
\question{Your advisor invented BEADS.  How do we know this thesis is not
just an exercise in making the advisor's algorithm look relevant?
Where is the objectivity?}

\begin{answerbox}
\textbf{Key points:}
\begin{itemize}[nosep]
  \item The thesis is arguably \emph{harder} on BEADS than an
        independent researcher would be: it identifies three failure
        modes of BEADS (sparsity violation, filter bandwidth mismatch,
        spatially uniform regularisation), demonstrates that the
        $\ell_1$ prior is fundamentally limited in dense-peak
        regimes, and shows that unrolling inherits these limitations.
  \item A thesis that just praised BEADS would be weaker.  The
        negative findings (leakage floor, over-regularisation,
        CG gradient collapse) are honest scientific reporting.
  \item BEADS was chosen because it has the cleanest structure for
        unrolling (explicit MM iterations, scalar parameters,
        joint peak+baseline estimation).  This is a technical
        motivation, not a promotional one.
  \item The thesis \emph{motivates moving beyond BEADS}
        (group sparsity, U-HQ penalties) as the primary future
        direction.
\end{itemize}
\textbf{Thesis references:} Sec.~2.2.1.5 (BEADS failure modes), Sec.~5.1.2 (leakage analysis), Sec.~5.7 (honest summary).
\end{answerbox}

% ── Q56 ─────────────────────────────────────────────────────────────────────
\question{Can you derive the soft-thresholding operator on the board right
now?  Show me that it is the proximal operator of the $\ell_1$
norm, and explain why it causes shrinkage bias.}

\begin{answerbox}
\textbf{Whiteboard derivation (MUST be prepared):}
\begin{enumerate}[nosep]
  \item The proximal operator: $\text{prox}_{\lambda\|\cdot\|_1}(\mathbf{v})
        = \arg\min_{\mathbf{x}} \frac{1}{2}\|\mathbf{x} - \mathbf{v}\|^2
        + \lambda\|\mathbf{x}\|_1$.
  \item The problem separates element-wise:
        $\min_{x_i} \frac{1}{2}(x_i - v_i)^2 + \lambda|x_i|$.
  \item Take subgradient, set to zero:
        $(x_i - v_i) + \lambda \cdot \text{sign}(x_i) = 0$.
  \item Three cases: $v_i > \lambda \Rightarrow x_i = v_i - \lambda$;
        $v_i < -\lambda \Rightarrow x_i = v_i + \lambda$;
        $|v_i| \le \lambda \Rightarrow x_i = 0$.
  \item This is $\mathcal{S}_\lambda(v_i)
        = \text{sign}(v_i)\max(|v_i| - \lambda, 0)$.
  \item \textbf{The bias}: for any non-zero true signal $x^* > 0$,
        $\mathcal{S}_\lambda(x^*) = x^* - \lambda < x^*$.  The
        estimate is \emph{always} shrunk by exactly $\lambda$.
        In dense-peak regions where many entries are non-zero,
        this shrinkage accumulates.
\end{enumerate}
\textbf{Thesis references:} Eq.~3.14 (asymmetric soft-thresholding), Sec.~3.2.4 (proximal step).
\end{answerbox}

% ── Q57 ─────────────────────────────────────────────────────────────────────
\question{Your masked baseline loss (Term~5) supervises the baseline only
where peaks are absent.  But at inference time, you do not know
where peaks are.  The model was trained with ground-truth peak
masks that it will never have in deployment.  Isn't the model
learning to rely on information it cannot access?}

\begin{answerbox}
\textbf{Key points:}
\begin{itemize}[nosep]
  \item The masked loss provides a \emph{training signal}, not
        an inference-time input.  The network does not receive the
        mask at test time---it receives only $\mathbf{y}$.
  \item The training objective teaches the network to produce
        baselines that are accurate in non-peak regions.  At inference,
        the network generalises this learned behaviour to unseen signals.
  \item This is standard supervised learning: the ground truth is used
        to compute the loss during training but is not available at
        test time.  All supervised losses ``rely on information the
        model cannot access at inference.''
  \item The concern would be valid if the mask were an \emph{input}
        to the network.  It is not---it only shapes the gradient.
  \item Potential issue: if the peak mask definition ($\gamma = 0.02$
        threshold) is very different between training and real data,
        the learned baseline behaviour may not transfer.
\end{itemize}
\textbf{Thesis references:} Sec.~3.7.3 (masked baseline loss), Eqs.~3.27--3.29.
\end{answerbox}

% ── Q58 ─────────────────────────────────────────────────────────────────────
\question{You claim ``interpretability'' as a key advantage.  But can you
actually \emph{diagnose} why LBEADS-NET fails on a specific signal
by inspecting the learned parameters?  Show me how.}

\begin{answerbox}
\textbf{Key points:}
\begin{itemize}[nosep]
  \item Partial yes: the per-layer parameters have physical meaning.
        If $\lambda_0^{(k)}$ is very large in early layers, one can
        diagnose ``early layers are over-thresholding small peaks.''
        If $r^{(k)}$ is low, ``insufficient asymmetric penalisation
        of negative artefacts.''
  \item More concretely: one can run the signal through each layer
        sequentially, inspecting the intermediate $\mathbf{x}^{(k)}$
        and $\mathbf{f}^{(k)}$ at each stage.  This ``layer-by-layer
        decomposition'' is impossible with a black-box U-Net.
  \item However, with $K{=}20$ layers and $5$ parameters each, the
        interaction effects are complex.  ``Interpretable'' does not
        mean ``simple to diagnose.''
  \item The emergent specialisation (Sec.~5.5) demonstrates that
        the parameters \emph{are} interpretable post-training.  But
        using them for per-signal failure diagnosis is an unsolved
        practical problem.
  \item Honest answer: ``more interpretable than a black box'' is
        correct, but ``fully diagnosable'' is an overstatement.
\end{itemize}
\textbf{Thesis references:} Sec.~5.5 (parameter analysis), Sec.~3.4.3 (layer chaining).
\end{answerbox}

% ── Q59 ─────────────────────────────────────────────────────────────────────
\question{You exploit dense $O(N^2)$ matrices when the orignal BEADS uses
$O(N)$ banded solvers.  An undergraduate could implement banded
matrix-vector products in PyTorch.  Why didn't you?}

\begin{answerbox}
\textbf{Key points:}
\begin{itemize}[nosep]
  \item Concede: banded matrix operations are straightforward in
        theory.  The challenge is integrating them cleanly with
        PyTorch's autograd for backpropagation.
  \item The highpass filter involves $\mathbf{H} = \mathbf{B}\mathbf{A}^{-1}$,
        where $\mathbf{A}$ is banded.  The forward pass requires
        solving $\mathbf{A}\mathbf{z} = \mathbf{v}$ (banded system solve).
        The \emph{backward} pass requires differentiating through
        this solve, which PyTorch does not natively support for
        banded systems.
  \item Options: (a)~use \texttt{torch.linalg.solve\_triangular}
        with LU-factored $\mathbf{A}$, (b)~write a custom autograd
        function for banded solves, (c)~use implicit differentiation.
  \item All of these require more engineering effort than converting
        to dense, which was the pragmatic choice given the thesis
        timeline.
  \item This is the single most important engineering improvement
        for practical deployment.
\end{itemize}
\textbf{Thesis references:} Sec.~3.2.2 (filter construction), Sec.~3.4 (dense implementation).
\end{answerbox}

% ── Q60 ─────────────────────────────────────────────────────────────────────
\question{ADMM-Net is the standard framework for unrolling constrained
optimisation problems.  Why did you not use ADMM?  The baseline
non-negativity constraint and the frequency separation are natural
ADMM constraints.}

\begin{answerbox}
\textbf{Key points:}
\begin{itemize}[nosep]
  \item ADMM decomposes the problem into sub-problems that handle
        different constraints/penalties separately, linked by dual
        variables (Lagrange multipliers).  This would naturally
        separate peak sparsity, baseline smoothness, and frequency
        separation into distinct ADMM updates.
  \item BEADS is formulated as an unconstrained MM problem, not as
        a constrained optimisation.  Converting to ADMM would require
        reformulating the cost function with explicit constraints,
        which changes the algorithm being unrolled.
  \item The thesis's goal is to unroll \emph{BEADS specifically},
        preserving its structure.  ADMM-based unrolling would be a
        different method.
  \item ADMM introduces additional hyperparameters (penalty parameter
        $\rho$, dual variable update rate) that would also need to be
        learned per layer.
  \item Concede: ADMM unrolling is a promising alternative approach
        that should be explored.
\end{itemize}
\textbf{Thesis references:} Sec.~2.3.4 (ADMM-Net survey), Sec.~3.2 (ISTA choice rationale).
\end{answerbox}

% ── Q61 ─────────────────────────────────────────────────────────────────────
\question{Your model is trained at $N{=}4096$ and tested at $N{=}4096$.
Real chromatograms can have $N{=}10^5$ to $10^6$ points.  What
happens when you apply this model to a different signal length?
Does it even work?}

\begin{answerbox}
\textbf{Key points:}
\begin{itemize}[nosep]
  \item \textbf{It does not work directly.}  The dense matrices
        $\mathbf{A}$, $\mathbf{B}$, $\mathbf{L}$ are precomputed at
        construction time for a fixed $N$.  A different signal length
        requires reconstructing these matrices (which changes their
        spectral properties) and re-registering the buffers.
  \item The 100 learned scalar parameters are \emph{independent of
        $N$}---they can be transferred to any signal length.  Only the
        filter matrices need to be recomputed.
  \item This is a strength of the architecture: the learned
        parameters are signal-length-agnostic.  The filter construction
        adapts automatically via \texttt{BAfilt(d, fc, N\_new)}.
  \item However, \textbf{this was not tested}.  The normalised cutoff
        $f_c{=}0.006$ has a different physical meaning at different $N$
        (the period $1/f_c \approx 167$ samples is fixed, so longer
        signals have proportionally more baseline variation relative
        to the filter's passband).
  \item With banded implementation, arbitrary $N$ would be trivial.
        With dense matrices, $N{=}10^5$ requires $\sim$80\,GB per
        matrix---infeasible.
\end{itemize}
\textbf{Thesis references:} Sec.~3.4 (architecture), Sec.~3.2.2 (filter parameterisation).
\end{answerbox}

% ── Q62 ─────────────────────────────────────────────────────────────────────
\question{You initialise all layers with identical parameters and let
training differentiate them.  But you only train for 30 epochs
(300 gradient steps).  With 100 parameters to differentiate in
300 steps, haven't the parameters barely moved from their
initialisation?  How much specialisation can actually emerge?}

\begin{answerbox}
\textbf{Key points:}
\begin{itemize}[nosep]
  \item The 30-epoch / 300-step budget is indeed very tight for
        100 parameters.  This is why the more complex loss
        configurations (Exp.~1.5) fail---they need more steps.
  \item However, the progressive experiments show meaningful
        learning: Peak MSE drops from 21.27 (classical) to 9.30
        (Exp.~1.3) in 300 steps.  The parameters have clearly moved.
  \item The parameter analysis in Sec.~5.5 uses a different model
        (6 layers, 100 epochs = 1000 steps) where specialisation
        is more pronounced.
  \item Concede: the 30-epoch budget was a practical constraint
        (training time on CPU), not an optimised choice.  Longer
        training would likely improve all configurations, especially
        the complex ones.
  \item This is precisely why the curriculum training (Sec.~3.6)
        exists: it provides a \emph{structured path} through parameter
        space that achieves more in fewer steps than random
        initialisation.
\end{itemize}
\textbf{Thesis references:} Table~4.2 (training config), Sec.~5.5 (parameter analysis), Sec.~3.6 (curriculum).
\end{answerbox}

% ── Q63 ─────────────────────────────────────────────────────────────────────
\question{Write the BEADS cost function on the board.  Now take the
gradient with respect to $\mathbf{x}$.  Now show me how one MM
iteration relates to one ISTA step.  Where does the approximation
enter?}

\begin{answerbox}
\textbf{Whiteboard derivation (MUST be prepared):}
\begin{enumerate}[nosep]
  \item BEADS cost: $J(\mathbf{x}) = \tfrac{1}{2}\|H(\mathbf{y} - \mathbf{x})\|^2
        + \lambda_0 \sum_n \theta(x_n) + \lambda_1 \sum_n \phi(\Delta x_n)
        + \lambda_2 \sum_n \phi(\Delta^2 x_n)$.
  \item Split into smooth + non-smooth:
        $J = \underbrace{J_{\text{smooth}}}_{\text{data fidelity + smoothness}}
        + \underbrace{J_{\text{non-smooth}}}_{\lambda_0 \sum \theta(x_n)}$.
  \item $\nabla J_{\text{smooth}} = -H^T H(\mathbf{y} - \mathbf{x})
        + \lambda_1 D_1^T \psi(D_1 \mathbf{x})
        + \lambda_2 D_2^T \psi(D_2 \mathbf{x})$
        where $\psi(z) = z/(|z| + \epsilon)$.
  \item \textbf{MM step}: majorise the non-quadratic terms, solve the
        resulting banded linear system exactly.  This gives
        \emph{exact} majoriser minimisation.
  \item \textbf{ISTA step}: take one gradient step on $J_{\text{smooth}}$:
        $\mathbf{x}_{\text{update}} = \mathbf{x}^{k-1} + \eta \nabla J_{\text{smooth}}$,
        then apply proximal operator for $\theta$:
        $\mathbf{x}^k = \mathcal{S}_{\text{asym}}(\mathbf{x}_{\text{update}};
        \lambda_0 \eta, r)$.
  \item \textbf{The approximation}: ISTA replaces exact MM
        minimisation with a single gradient step + proximal.  It is
        less accurate per iteration but \emph{much simpler} to
        differentiate.
\end{enumerate}
\textbf{Thesis references:} Eq.~2.1 (BEADS cost), Algorithm~1 (MM), Algorithm~2 (ISTA layer), Sec.~3.3.
\end{answerbox}

% ── Q64 ─────────────────────────────────────────────────────────────────────
\question{Your v5 model uses clamped exponential parameterisation, but
Section~3.2.1 says ``No clamping is applied---the exponential
function alone guarantees positivity.''  Which is it?  Does the
6-layer parameter analysis (Table~5.5) show clamped parameters
(marked with $*$)?  These statements contradict each other.}

\begin{answerbox}
\textbf{Key points:}
\begin{itemize}[nosep]
  \item This is a genuine inconsistency in the thesis.
  \item Section~3.2.1 describes the v5 architecture (the final model)
        where parameters use unclamped exponential parameterisation:
        $\lambda = \exp(\theta)$, no bounds.
  \item Table~5.5 shows parameters from a \emph{different model}
        (6-layer, $N{=}1024$, curriculum-trained) that uses clamped
        parameterisation (parameters marked with $*$ hit clamp
        boundaries).
  \item The two models have different parameterisation schemes.  This
        is not made clear in the thesis and creates confusion.
  \item Fix: add a note in Sec.~5.5 that the 6-layer model uses a
        different (earlier) parameterisation with clamps, and clarify
        that the main 20-layer model uses unclamped exponential.
\end{itemize}
\textbf{Thesis references:} Sec.~3.2.1 (unclamped), Table~5.5 (clamped $*$ markers).
\end{answerbox}

% ── Q65 ─────────────────────────────────────────────────────────────────────
\question{You claim the per-layer independence enables ``emergent parameter
specialisation.''  But with 100 parameters trained on 160 signals
for 30 epochs, how do you distinguish ``specialisation'' from
``overfitting to the training set''?  What would the parameters
look like on a different random seed?}

\begin{answerbox}
\textbf{Key points:}
\begin{itemize}[nosep]
  \item Valid concern.  The monotonic increase in $\lambda_0$ across
        layers could be a genuine pattern or a training artefact.
  \item To distinguish: train multiple models with different random
        seeds and check whether the same qualitative pattern
        (increasing sparsity, non-monotonic asymmetry, stable step
        size) appears.  This was not done.
  \item Evidence against overfitting: the test-set metrics improve
        consistently across the progressive experiments, suggesting
        the learned parameters generalise.  If the model were
        overfitting, we would expect test performance to diverge from
        training performance.
  \item The specialisation pattern is also \emph{physically sensible}:
        it matches the behaviour of iterative thresholding algorithms
        (coarse-to-fine refinement).  Overfitting would produce
        arbitrary or noisy parameter values.
  \item Concede: multi-seed experiments are needed to confirm the
        robustness of the specialisation pattern.
\end{itemize}
\textbf{Thesis references:} Sec.~5.5 (parameter analysis), Table~5.5 (learned values).
\end{answerbox}


% ============================================================================
%  KNOWN INCONSISTENCIES TO FIX
% ============================================================================
\newpage
\section{Known Inconsistencies \& Issues to Fix Before Defense}

\begin{warnbox}[CRITICAL: Fix Before Defense]
The following issues have been identified in the current thesis draft.
Each should be resolved before the defense to avoid uncomfortable questions.
\end{warnbox}

\vspace{0.5em}

\renewcommand{\arraystretch}{1.4}
\begin{table}[h!]
\centering
\small
\begin{tabular}{@{} c p{5.8cm} p{3cm} p{3.5cm} @{}}
\toprule
\textbf{\#} & \textbf{Issue} & \textbf{Location} & \textbf{Fix} \\
\midrule
1 & Abstract is Lorem Ipsum placeholder text &
    Abstract page &
    Write the real abstract \\
2 & Conclusions chapter is Lorem Ipsum placeholder &
    Chapter 6 &
    Write real conclusions \\
3 & Dedication has ``YOUR DEDICATION (1-2 LINES)'' placeholder &
    Dedication page &
    Write dedication or remove \\
4 & Broken cross-references showing ``Chapter~??'' &
    pp.\ 43, 56, 68 &
    Fix \verb|\label/\ref| pairs \\
5 & University ID and NetID are placeholders (``N\#'', ``NetID'') &
    Title/signature pages &
    Insert real ID numbers \\
6 & Ablation tables 5.3 and 5.4 have all dashes &
    Tables 5.3--5.4 &
    Run experiments or remove \\
7 & Sections 5.3 (Qualitative Examples) and 5.4 (Training Dynamics) are placeholder &
    Sec.\ 5.3--5.4 &
    Write or mark as future work \\
8 & Advisor title inconsistency: ``Assistant Prof.'' vs ``Professor'' &
    Abstract vs.\ approval page &
    Verify correct title, unify \\
9 & Loss term count: ``11-term'' in Intro vs ``twelve'' in Sec.\ 3.7 &
    Sec.\ 1.2 vs Sec.\ 3.7 &
    Correct Intro to ``twelve-term'' \\
10 & Peak model mismatch: EMG in Sec.\ 2.1.4 vs bi-Gaussian in Sec.\ 4.1.1 &
     Sec.\ 2.1.4 vs 4.1.1 &
     Add justification in Sec.\ 4.1.1 \\
\bottomrule
\end{tabular}
\caption{Known issues in the current thesis draft, ordered by severity.}
\label{tab:issues}
\end{table}

\vspace{1em}

\begin{warnbox}[Priority Actions]
\begin{enumerate}[nosep]
  \item[\textbf{P0}] Items 1, 2, 5: These will be immediately noticed by any
        committee member opening the document.
  \item[\textbf{P1}] Items 3, 4, 8, 9: Visible quality issues that suggest
        the document was not carefully reviewed.
  \item[\textbf{P2}] Items 6, 7, 10: Content gaps that may or may not be
        raised depending on committee member interests.
\end{enumerate}
\end{warnbox}


% ============================================================================
%  GENERAL DEFENSE ADVICE
% ============================================================================
\newpage
\section*{General Defense Advice}
\addcontentsline{toc}{section}{General Defense Advice}

\begin{tcolorbox}[
  enhanced,
  colback=blue!3!white,
  colframe=SectionBlue,
  fonttitle=\bfseries\sffamily,
  title={Guiding Principles for the Defense},
  coltitle=white,
  attach boxed title to top center={yshift=-2mm},
  boxed title style={colback=SectionBlue, colframe=SectionBlue, sharp corners, size=small},
  left=6pt, right=6pt, top=6pt, bottom=6pt
]

\begin{enumerate}[leftmargin=*, itemsep=8pt]

\item \textbf{Own the negative results.}
The BLI plateau at 0.59 and the over-regularisation finding (Exp.~1.5)
are \emph{contributions}, not failures.  Present them as deliberate
investigations.  ``We tried X, and it did not work because of Y'' is a
perfectly valid scientific finding.

\item \textbf{Know what you do not know.}
When a question targets a gap (e.g., ``why no arPLS comparison?''),
say: ``That is a limitation of this work.  Here is what I would do
if I had more time.''  Do not bluff.

\item \textbf{Redirect to the architecture.}
Many attack questions (area error, BLI, peak correlation) have the same
root cause: $\ell_1$ shrinkage bias.  Learn the shrinkage bias argument
cold and redirect to it.

\item \textbf{Prepare the whiteboard.}
Be ready to derive:
(a)~the BEADS cost function (Eq.~2.24),
(b)~the ISTA proximal step (Eq.~3.9),
(c)~the soft-thresholding operator and its bias,
(d)~why untied parameters break MM convergence.

\item \textbf{Credit the advisor.}
Professor Selesnick co-authored BEADS.  When discussing the algorithm,
frame the work as a natural evolution of the group's research programme.

\item \textbf{Time management.}
Aim for 25--30 minutes of presentation, leaving 30--45 minutes for
questions.  For each question, answer in 60--90 seconds.  If the
committee wants more detail, they will ask.

\item \textbf{Do not read from this document during the defense.}
Internalise the key points.  The goal is fluency, not recitation.

\end{enumerate}
\end{tcolorbox}

\vfill
\begin{center}
\small\itshape
Prepared April 2026.\\
Good luck, Saad.
\end{center}


% ============================================================================
%  GLOSSARY
% ============================================================================
\newpage
\section*{Glossary of Key Terms}
\addcontentsline{toc}{section}{Glossary of Key Terms}

\begin{tcolorbox}[
  enhanced,
  colback=blue!3!white,
  colframe=SectionBlue,
  fonttitle=\bfseries\sffamily,
  title={Terms Defined in the Context of This Thesis},
  coltitle=white,
  attach boxed title to top center={yshift=-2mm},
  boxed title style={colback=SectionBlue, colframe=SectionBlue, sharp corners, size=small},
  left=6pt, right=6pt, top=6pt, bottom=6pt
]
{\small This glossary defines terms as they are used specifically within the
LBEADS-NET thesis.  Definitions are scoped to the chromatographic
baseline-correction context and may differ from their broader usage in
other fields.}
\end{tcolorbox}

\vspace{1em}

\newcommand{\glossterm}[1]{\par\vspace{0.6em}\noindent\textbf{\textsf{\textcolor{SectionBlue}{#1}}}\par\nopagebreak\vspace{0.15em}}

\glossterm{$\ell_1$ Leakage Limit}
The empirically observed floor on the Baseline Leakage Index
(BLI~$\approx 0.59$) that persists across \emph{all} LBEADS-NET loss
configurations (Experiments~1.1--1.5).  The limit arises because the
$\ell_1$-based soft-thresholding operator at the core of every ISTA layer
shrinks \emph{every} non-zero signal value by exactly~$\lambda$
(the effective threshold).  In dense-peak regions where the sparsity
assumption is violated, this shrinkage accumulates: the total missing
peak energy is proportional to $\lambda \times (\text{number of non-zero
samples})$, and this energy is absorbed into the baseline estimate via the
low-pass filter.  Because the shrinkage is a structural property of the
proximal operator---not a tunable weight---no amount of loss engineering
(anti-leakage penalties, asymmetric baseline loss, envelope constraints,
etc.)\ can push the BLI below this floor.  The limit generalises to
\emph{any} unrolled network that inherits $\ell_1$-based sparsity priors,
as independently confirmed by Gharbi et~al.\ (2024) in their work on
unrolled sparse-recovery architectures.
\hfill\textit{Thesis refs: Sec.~5.1.2, Table~5.2, Sec.~1.3 (Contribution~3)}


\end{document}
